\chapter{Metodologia}\label{cap:metodologia}

% Manual (trabalho prático): ferramentas usadas, critérios de escolha e por que
% não outras; os 4 diagramas UML (casos de uso, classes, sequência e atividades)
% e o MER, cada um descrito; público-alvo definido. Texto no passado e na voz
% passiva ("foi realizado", "optou-se por").

\section{Classificação da pesquisa}

\pendente{Pesquisa aplicada, de abordagem quantitativa e procedimento experimental. Justificar cada classificação.}

\section{Público-alvo}

\pendente{Produtores de tomate, sobretudo pequenos e familiares, que precisam decidir o momento da colheita no campo, muitas vezes sem internet.}

\section{Etapas do trabalho}

\pendente{Fluxograma das etapas: escolha do dataset, definição das classes, treinamento, desenvolvimento do aplicativo, testes e avaliação.}

\section{Ferramentas e justificativas}

\pendente{Para cada ferramenta (Flutter, pacote Ultralytics e YOLO26n-cls, LiteRT e Core ML, SQLite, Google Colab, Overleaf, GitHub), explicar o critério da escolha e por que outras opções não foram usadas.}

\section{Conjunto de dados e classes}\label{sec:dados}

\pendente{Dataset público usado (origem, licença, número de imagens), mapeamento das classes originais para verde, de vez, maduro, passado e outro, e como foram gerados os recortes. Conjunto de teste próprio: tomates comprados verdes e fotografados diariamente até amadurecerem, com 2 ou 3 celulares e variação de luz. Separação entre treino e teste por fruto ou lote, nunca por foto.}

\section{Treinamento do modelo}

\pendente{Versão fixada do pacote, pesos pré-treinados, tamanho de imagem, número de épocas, hardware e semente. Tudo deve ser reprodutível a partir do notebook do repositório.}

\section{Requisitos}

\begin{table}[htb]
  \ABNTEXfontereduzida
  \caption{Requisitos funcionais e não funcionais do aplicativo}
  \label{tab:requisitos}
  \begin{tabularx}{\textwidth}{lX}
    \hline
    \textbf{Código} & \textbf{Requisito} \\
    \hline
    RF01 & Capturar uma foto pela câmera ou escolher uma imagem da galeria. \\
    RF02 & Classificar a maturação do tomate no próprio dispositivo, sem conexão com a internet. \\
    RF03 & Exibir a classe, a confiança e a recomendação de colher ou esperar. \\
    RF04 & Alertar o usuário quando a foto for inadequada ou a confiança for baixa. \\
    RF05 & Armazenar cada análise no histórico do dispositivo. \\
    RF06 & Consultar o histórico de análises. \\
    RNF01 & Funcionar integralmente sem conexão com a internet. \\
    RNF02 & Executar em Android 7.0 ou superior. \\
    RNF03 & Responder em até \pendente{META DE TEMPO} por análise. \\
    RNF04 & Não enviar fotos nem dados de localização para servidores. \\
    RNF05 & Manter a interface legível sob sol forte, com alto contraste e botões grandes. \\
    \hline
  \end{tabularx}
  \legend{Fonte: elaborada pelos autores.}
\end{table}

\section{Modelagem do sistema}

\subsection{Diagrama de casos de uso}
\pendente{Figura e descrição: fotografar fruto, analisar maturação, ver recomendação e consultar histórico. Ator: produtor.}

\subsection{Diagrama de classes}
\pendente{Figura e descrição: Analise, Cultura, ClasseMaturacao, RegraDecisao, ClassificadorLocal e RepositorioAnalises.}

\subsection{Diagrama de sequência}
\pendente{Figura e descrição: foto, pré-processamento, classificador, regra de decisão, tela de resultado e histórico.}

\subsection{Diagrama de atividades}
\pendente{Figura e descrição, com as decisões: a foto está boa? a confiança é suficiente? a recomendação é colher ou esperar?}

\section{Modelo entidade-relacionamento}

\pendente{Figura e descrição: Cultura (1:N) ClasseMaturacao (1:1) RegraDecisao; Analise (N:1) ClasseMaturacao, com data, caminho da foto no dispositivo, confiança e versão do modelo.}

\section{Plano de testes e de avaliação}

\pendente{Casos de teste derivados dos requisitos; métricas do modelo (acurácia, precisão, recall, F1 por classe e matriz de confusão) no conjunto de teste próprio; tempo de resposta medido em 2 ou 3 celulares; funcionamento em modo avião.}

\section{Declaração de uso de inteligência artificial}

Durante a elaboração deste trabalho, foram utilizadas ferramentas de inteligência artificial generativa como apoio à pesquisa bibliográfica, à escrita e revisão do texto e ao desenvolvimento do código. Todo o conteúdo produzido com esse apoio foi revisado pelos autores, as referências citadas foram conferidas nas fontes originais e todos os resultados apresentados vêm de execuções reais do modelo e do aplicativo. \pendente{Listar as ferramentas usadas, com nome e fabricante.}
